\section{Test og kvalitetssikring}

Der har gennem hele projektforløbet været stort fokus på en agil arbejdsproces med henblik på hurtigt at kunne levere en Minimal Viable Product (MVP), som kunne anvendes i det danske kajakpolo-miljø.
Som følge heraf har der ikke været et stort fokus på etablering af omfattende automatiserede tests eller CI/CD-pipelines til automatisk eksekvering af test ved hver kodeændring.
I stedet er kvalitetssikringen primært blevet udført gennem manuelle tests i forbindelse med både implementering af nye funktioner og ændringer i eksisterende kode.
Dette vurderes acceptabelt i projektets nuværende fase, hvor systemets kompleksitet fortsat er relativt begrænset.
Der har dog været krav om, at alle kodeændringer blev reviewet af et andet gruppemedlem, med henblik på at sikre kvalitet og reducere risikoen for fejl.

Denne tilgang har gjort det muligt at accelerere i udviklingsprocessen, implementere nye features hurtigere, samt håndtere fejlrettelser mere effektivt.
Samtidig har der været et kontinuerligt fokus på at sikre en høj kvalitet i produktet, selv uden fuld automatisering af testprocessen.

Der er dog implementeret en automatiseret test af algoritmen, som genererer turneringsplaner.
Dette skyldes, at den genererede datamængde kan være så omfattende, at manuel validering ville være både tidskrævende og upraktisk.
Turneringsgeneratoren skal overholde en række centrale regler, eksempelvis at et hold ikke må spille to kampe samtidigt, hvis der er flere baner til rådighed,
samt at et hold maksimalt må spille to kampe i træk uden pause.

Disse regler sikrer, at den genererede turneringsplan kan opfylder de praktiske krav til afvikling af turneringen uden omfattende manuel justering.

Kvalitet er en central faktor for systemets anvendelighed, og derfor har der været et betydeligt fokus på kvalitetssikring, selvom automatiserede tests ikke er fuldt integreret i udviklingspipelinen.

\subsection{Turneringstest}
For at teste turneringsgeneratoren, er der testframework \texttt{Jest} anvendt (\cite{Jest}).

I filen \texttt{tournament.service.test.ts} findes testkoden for den del af systemet, der vurderes at have størst betydning for produktets funktionalitet, nemlig algoritmen der genererer turneringsplaner.

Formålet med testen er at sikre, at den genererede turneringsplan overholder de centrale forretningsregler for afvikling af en kajakpoloturnering.
Denne del af produktet, er hvor input varierer højest, og er samtidig den mest efterspurgte funktion fra vores Product Owner.
Derfor er det nødvendigt at sikre en vis robusthed i genereringen af en turneringsplan.

Testen verificerer bl.a. følgende krav:

\begin{itemize}
    \item Et hold må ikke være planlagt til to kampe på samme tidspunkt
    \item Et hold må maksimalt spille to kampe i træk uden en pause
    \item Alle hold skal møde hinanden to gange
    \item Turneringsplanen skal indeholde X antal kampe hvor $X = n(n-1)$ og \textit{n} er antal hold
\end{itemize}

I stedet for at sammenligne med ét fast forventet output, valideres turneringsplanen ud fra disse constraints.
Dette skyldes, at flere forskellige turneringsplaner kan være valide, så længe de opstillede regler overholdes.

Denne test fungerer samtidig som en regressionstest, der gør det muligt hurtigt at verificere, at fremtidige ændringer i algoritmen ikke introducerer fejl i planlægningslogikken.